\chapter{Extending the portal}

This chapter is the load-bearing summary of the three slide decks under
\reposrc{docs/publishing/slides/}: how to add a new game, a new role, a
new federated route, or a new admin view. The slides have screenshots
and click-by-click walkthroughs; this chapter is the canonical
\emph{contract} those walkthroughs follow.

\section{Adding a new game}

There are two paths.

\subsection{Path A — upload at runtime}

Every supported game ships an editor plugin
(\snapfile{addons/snapdragon_export/}) under
\reposrc{games/three-wizards-in-a-trenchcoat/} that produces a deploy-ready
zip with a v2 \texttt{manifest.json}, a \texttt{register.sql} insert
script, and the \texttt{bundle/} contents:

\begin{enumerate}
  \item In the Godot editor, \emph{Project~$\rightarrow$~Tools~$\rightarrow$~Snapdragon~$\rightarrow$~Export\dots}.
  \item Fill in slug, display name, description, scene type. The plugin
        defaults the federation \texttt{remote\_entry\_url} to
        \texttt{/\_remotes/games-<slug>/remoteEntry.json}; if you don't
        plan to ship a per-game frontend remote, change it to
        \texttt{/\_remotes/wasm-fragment/remoteEntry.json} so the host
        loads the bundle through the shared fragment.
  \item Click \emph{Export}. The plugin writes a zip to your chosen
        path.
  \item Hit \texttt{POST /api/admin/games/upload} with the zip — the
        admin UI's \emph{Upload game bundle} card does this for you.
        The install transaction validates, hashes, atomically moves the
        bundle, and runs \texttt{register.sql} which inserts into
        \texttt{games}, \texttt{federation\_modules}, and
        \texttt{game\_access} all in one Postgres transaction.
\end{enumerate}

The host shell discovers the new game on the next federation manifest
fetch (currently every 60s; see \snapfile{FederationService}). No host
redeploy is required.

\subsection{Path B — commit the game}

For a game you want to ship in the repository:

\begin{enumerate}
  \item Drop the Godot project under
        \snapfile{games/<game-dir>/project.godot}, with its own
        \snapfile{export_presets.cfg} and
        \snapfile{addons/snapdragon_export/}.
  \item Add the directory to \texttt{GAME\_PROJECTS} and define
        \texttt{PRESETS\_<dir>} in the root \snapfile{Makefile}.
  \item Add a Liquibase changeset that inserts the \texttt{games} +
        \texttt{federation\_modules} + \texttt{game\_access} rows. See
        \reposrc{backend/src/main/resources/db/changelog/017-snapdragon-catalog.xml}
        for the canonical pattern.
  \item Optionally add a \texttt{games-<slug>} service to
        \snapfile{docker-compose.yml} if the game ships its own
        Module Federation remote (separate Angular project under
        \snapfile{frontend/projects/games-<slug>/}).
  \item Run \texttt{make build} and \texttt{docker compose up -d}.
\end{enumerate}

\section{Adding a new federated route}

\begin{enumerate}
  \item Build the route's component as a standalone Angular project
        under \snapfile{frontend/projects/<remote-name>/}, with its own
        \snapfile{federation.config.js} extending
        \snapfile{frontend/federation-base.js}.
  \item Expose the entry component in
        \snapfile{frontend/projects/<remote-name>/federation.config.js}'s
        \texttt{exposes} map.
  \item Add a row to \texttt{federation\_modules} (Liquibase changeset
        recommended). Set \texttt{kind} to \texttt{remote-shell} for a
        full route component.
  \item In \snapfile{frontend/src/app/app.routes.ts}, plug
        \texttt{loadRemote('<remote-name>', './<ExposedSymbol>',
        '<ExposedSymbol>')} into a \texttt{loadComponent} entry.
  \item Add an nginx \texttt{location /\_remotes/<remote-name>/} block
        in \snapfile{frontend/nginx.conf} pointing at the per-remote
        compose service.
  \item Add the per-remote service to \snapfile{docker-compose.yml}.
\end{enumerate}

\section{Adding a new admin view}

The admin remote (\reposrc{frontend/projects/admin/}) exposes any
top-level standalone component listed in its \texttt{exposes} map. To
add a view:

\begin{enumerate}
  \item Add a standalone component under
        \snapfile{frontend/projects/admin/src/app/}.
  \item Add it to the \texttt{exposes} map in
        \snapfile{frontend/projects/admin/federation.config.js}.
  \item Wire a route in the host's
        \snapfile{frontend/src/app/app.routes.ts} pointing at the
        exposed name.
\end{enumerate}

DTOs cannot be shared across the federation boundary at runtime
(TypeScript types erase). Redeclare interfaces inside the remote — see
the existing pattern in
\reposrc{frontend/projects/admin/src/app/upload-remote.component.ts}.

\section{Adding a new role}

Casdoor is the source of truth for the role catalog. The portal
enforces visibility through two layers:

\begin{enumerate}
  \item Create the role in Casdoor (\texttt{ROLE\_PLAY\_<GAME>} for game
        access, \texttt{ROLE\_ADMIN} for admin). The
        \snapfile{scripts/seed-roles.sh} script creates the per-game
        \texttt{ROLE\_PLAY\_*} roles idempotently; run it after
        deploying a new game.
  \item Insert a \texttt{game\_access} row with
        \texttt{(scope='role', value='ROLE\_PLAY\_<GAME>')} via the
        admin UI's \emph{Per-game access} card or by direct SQL.
\end{enumerate}

The admin dashboard's \emph{Grant role} dropdown reads the live Casdoor
catalog through \texttt{GET /api/admin/roles}; new roles appear as soon
as Casdoor's catalog has them.

\section{Adding a new SSE channel}

\begin{enumerate}
  \item Add a \texttt{Sinks.Many} field to \snapfile{GameStateService}
        with \texttt{autoCancel=false} (see the warning in
        chapter~\ref{ch:datamodel}).
  \item Add a \texttt{ServerSentEvent}-builder method that wraps the
        sink as a typed \texttt{Flux} with an \texttt{event:} field.
  \item Merge it into the \texttt{Flux.merge(...)} call inside
        \snapfile{GameStreamController} alongside the existing four
        channels and the keepalive.
  \item Register a per-type listener on the frontend by calling
        \texttt{eventSource.addEventListener('<event-name>', cb)} from
        \snapfile{GameStateService} or wherever the new channel is
        consumed.
\end{enumerate}

\section{Adding a new database column}

\begin{enumerate}
  \item Write a Liquibase changeset under
        \snapfile{backend/src/main/resources/db/changelog/} with the
        next sequential number. Postgres-only — Cassandra schemas live
        in \snapfile{infra/cassandra/cassandra-init.cql}.
  \item Include the changeset in \snapfile{db.changelog-master.xml}.
  \item Add the field to the matching R2DBC entity. The R2DBC mapping
        is reflective; \snapfile{R2dbcConvertersConfig} only needs new
        entries when adding a new entity, not a new column.
  \item For GraalVM native-image builds, run \texttt{make
        build-backend-native} and verify the reflection hints still
        resolve. If a new entity, regenerate the hints with the
        \texttt{native-image-agent} per the
        \reposrc{README.md} \emph{GraalVM native-image build} section.
\end{enumerate}
